\begin{table}[H]
\centering\footnotesize
\caption{AP: ba cách mô hình hóa mùa vụ, có và không có độ trôi, AIC và RMSE trên tập kiểm định 1959--1960}\label{tab:tab_b12_10_seasonality_models}
\begin{adjustbox}{max width=\linewidth}
\begin{tabular}{|>{\raggedright\arraybackslash}p{\dimexpr 0.315\linewidth-2\tabcolsep\relax}|r|r|r|r|}
\toprule
Mô hình & k & n dòng & AIC & RMSE giữ lại \\
\midrule
airline & 3 & 107 & -389.0 & 43.2 \\[3pt]
11 biến giả + $\nabla$ + MA(1) & 13 & 119 & -424.3 & 85.5 \\[3pt]
11 biến giả + $\nabla$ + MA(1) + trôi & 14 & 119 & -434.0 & 33.1 \\[3pt]
Fourier K = 2 & 6 & 119 & -311.7 & 74.8 \\[3pt]
Fourier K = 2 + trôi & 7 & 119 & -330.5 & 23.6 \\[3pt]
Fourier K = 3 & 8 & 119 & -330.0 & 77.9 \\[3pt]
Fourier K = 3 + trôi & 9 & 119 & -347.5 & 21.9 \\[3pt]
Fourier K = 4 & 10 & 119 & -384.6 & 82.0 \\[3pt]
Fourier K = 4 + trôi & 11 & 119 & -398.0 & 27.6 \\[3pt]
\bottomrule
\end{tabular}
\end{adjustbox}\par\vspace{2pt}{\footnotesize\textit{Ghi chú:} n dòng: số quan sát còn lại sau sai phân, như Ví dụ 4.2 của Bài 12. AIC chỉ so sánh trực tiếp được giữa các mô hình cùng cách sai phân. Dòng ``+ trôi'' thêm hằng số vào phương trình của $\nabla y$, tức một độ trôi; Ví dụ 4.2 dùng bản không có độ trôi.}
\end{table}
